На сегодняшний день практически любое приложение использует базу данных для хранения и обработки
информации. Неудивительно, что существует множество подходов к организации системы хранения данных.
Различные решения могут отличаться по модели представления информации, способам обеспечения
целостности, масштабируемости и производительности. Выбор конкретного варианта зачастую является
неочевидной задачей, поскольку преимущества одной модели хранения данных в одних сценариях могут
сопровождаться ухудшением характеристик в других. Например, высокая эффективность при записи данных
не всегда сочетается с удобством и скоростью их извлечения, а гибкость структуры хранения может
достигаться ценой увеличения сложности поддержки целостности всей системы данных.

Для каждого приложения необходимо подобрать модель хранения данных, которая будет соответствовать
его требованиям и обеспечивать оптимальный баланс между производительностью, удобством
использования и поддержкой. Тем не менее часто можно руководствоваться общими рекомендациями.
Так, например, реляционные базы данных с нормализованной структурой хранения данных традиционно
считаются более подходящими для приложений, требующих строгой целостности данных и сложных связей
между сущностями, а документно-ориентированные базы данных, допускающие более гибкую структуру
хранения данных, могут лучше подходить для приложений, где структура данных может часто меняться.

Однако, чтобы гарантировать преимущества и недостатки различных моделей хранения данных в
конкретной предметной области, надежнее всего сравнить итоговую производительность в реализациях
на разных моделях или даже СУБД.

Примером такой "неочевидной" предметной области может служить система видеонаблюдения,
включающая большое количество камер. Они формируют большой поток информации в реальном времени.
Причем данные, которые необходимо хранить, могут быть разнородными и также часто меняться,
к примеру метаданные о событиях и объектах, обнаруженных на видео, или информацию о состоянии
камер и их конфигурации. При эксплуатации подобных систем не ясно, какая модель хранения данных
будет более эффективной в таких условиях. Для выбора конкретного варианта следует провести
исследование особенностей хранения данных системы видеонаблюдения, снять основные показатели
производительности и сравнить их для разных моделей хранения данных, чтобы сделать выбор в пользу
одной из них.

Целью данной работы является сравнительный анализ нормализованной и документно-ориентированной
моделей хранения данных системы видеонаблюдения на примере СУБД PostgreSQL и MongoDB.
Соответственно, необходимо четко определить критерии сравнения и сценарии использования. После
чего провести нагрузочное тестирование обеих моделей хранения данных, оценить итоговую
производительность и сделать выводы о том, какая модель хранения данных более эффективна для
системы видеонаблюдения.

